\documentclass[12pt]{article}
\usepackage[a3paper, margin=1in]{geometry}
\usepackage{amsmath, amssymb, amsthm, ulem, booktabs}

\begin{document}

\section*{Domácí úkol 10.2}

Vypracoval: Daniel \textit{``Randál"} Ransdorf \hfill Podpis: \rule{4cm}{0.4pt}

\begin{enumerate}
  
  \item Víme, že pro použití $[f]^B_B$ v předpisu $f$ musíme nejprve vektory
    převést do báze $B$, pak aplikovat matici $[f]^B_B$, pak převést zpět do kanonické báze.
    Pak platí:

    \begin{align*}
      A &= [id]^B_K \cdot [f]^B_B \cdot [id]^K_B \quad (*) \\
      A \cdot [id]^B_K  &= [id]^B_K \cdot [f]^B_B \cdot [id]^K_B \cdot [id]^B_K  \\
      A \cdot [id]^B_K  &= [id]^B_K \cdot [f]^B_B \\
      [id]^B_K = \left(\begin{array}{cc}
        a&b\\c&d
      \end{array}\right) \quad \Longrightarrow \quad
      \left(\begin{array}{cc}
        1&1\\1&0
      \end{array}\right) \left(\begin{array}{cc}
        a&b\\c&d
      \end{array}\right) &= \left(\begin{array}{cc}
        a&b\\c&d
      \end{array}\right) \left(\begin{array}{cc}
        3&0\\3&3
      \end{array}\right) \\
      \left(\begin{array}{cc}
        a+c & b+d \\
        a & b
      \end{array}\right) &= \left(\begin{array}{cc}
        3a+3b & 3b \\
        3c+3d & 3d
      \end{array}\right) \\
      \left(\begin{array}{cc}
        3a+2b+c & 3b+d \\
        a+2c+2d & b+2d
      \end{array}\right) &= \left(\begin{array}{cc}
        0&0\\0&0
      \end{array}\right) \\
      \left(\begin{array}{c}
        3a+2b+c \\ 3b+d \\
        a+2c+2d \\ b+2d
      \end{array}\right) &= \left(\begin{array}{c}
        0\\0\\0\\0
      \end{array}\right) \\
      \left(\begin{array}{cccc}
        3 & 2 & 1 & 0 \\
        0 & 3 & 0 & 1 \\
        1 & 0 & 2 & 2 \\
        0 & 1 & 0 & 2
      \end{array}\right) \left(\begin{array}{c}
        a\\b\\c\\d
      \end{array}\right) &= \left(\begin{array}{c}
        0\\0\\0\\0
      \end{array}\right) \\
      \left(\begin{array}{cccc|c}
        3 & 2 & 1 & 0 & 0 \\
        0 & 3 & 0 & 1 & 0 \\
        1 & 0 & 2 & 2 & 0 \\
        0 & 1 & 0 & 2 & 0
      \end{array}\right)
      \sim
      \left(\begin{array}{cccc|c}
        0 & 2 & 0 & 4 & 0 \\
        0 & 3 & 0 & 1 & 0 \\
        1 & 0 & 2 & 2 & 0 \\
        0 & 1 & 0 & 2 & 0
      \end{array}\right)
      \sim
      \left(\begin{array}{cccc|c}
        1 & 0 & 2 & 2 & 0 \\
        0 & 1 & 0 & 2 & 0 \\
        0 & 0 & 0 & 0 & 0 \\
        0 & 0 & 0 & 0 & 0 \\
      \end{array}\right)
      &\sim
      \left(\begin{array}{cccc|c}
        1 & 0 & 2 & 2 & 0 \\
        0 & 1 & 0 & 2 & 0 \\
        0 & 0 & 0 & 0 & 0 \\
        0 & 0 & 0 & 0 & 0 \\ \hline
        0 & 0 & 1 & 0 & s \\
        0 & 0 & 0 & 1 & t
      \end{array}\right)
      \sim
      \left(\begin{array}{cccc|c}
        1 & 0 & 0 & 0 & 3s+3t \\
        0 & 1 & 0 & 0 & 3t \\
        0 & 0 & 0 & 0 & 0 \\
        0 & 0 & 0 & 0 & 0 \\ \hline
        0 & 0 & 1 & 0 & s \\
        0 & 0 & 0 & 1 & t
      \end{array}\right) \\
      [id]^B_K = \left(\begin{array}{cc}
        a&b\\c&d
      \end{array}\right) &= \left(\begin{array}{cc}
        3s+3t & 3t \\ s & t
      \end{array}\right) \\
      [id]^B_K &= s\left(\begin{array}{cc}
        3 & 0 \\ 1 & 0
      \end{array}\right) + t\left(\begin{array}{cc}
        3 & 3 \\ 0 & 1
      \end{array}\right) \\
      [id]^B_K \text{ musí být regulární, zvolme}&\text{ tedy }s=0,t=1 \\
      [id]^B_K &= \left(\begin{array}{cc}
        3&3 \\ 0&1
      \end{array}\right)
    \end{align*} 

    Z matice přechodu do kanonické báze se dá báze ``přečíst" ... nalezli jsme bázi \underline{\underline{$B=\left((3,0)^T,(3,1)^T\right)$}} prostoru $\mathbb{Z}^2_5$, která splňuje zadání.

    \bigskip

    \footnotesize
    \textbf{Kontrola}: Nalezněme inverz k $[id]^B_K$ a ověřme (*)

    \[
      [id]^K_B = \left(\begin{array}{cc}
        3&3 \\ 0&1
      \end{array}\right)^{-1} = \left(\begin{array}{cc}
        2&4 \\ 0&1
      \end{array}\right)
    \]
    \[
      \left(\begin{array}{cc}
        1&1 \\ 1&0
      \end{array}\right) \overset{(*)}{=} \left(\begin{array}{cc}
        3&3 \\ 0&1
      \end{array}\right) \left(\begin{array}{cc}
        3&0 \\ 3&3
      \end{array}\right) \left(\begin{array}{cc}
        2&4 \\ 0&1
      \end{array}\right) =
      \left(\begin{array}{cc}
        3&3 \\ 0&1
      \end{array}\right) \left(\begin{array}{cc}
        1&2 \\ 1&0
      \end{array}\right) = \left(\begin{array}{cc}
        1&1 \\ 1&0
      \end{array}\right) \quad \checkmark
    \]

    

\end{enumerate}

\end{document}
